\section{结论}

\subsection{问题一：温度响应与时间变化}

在附件1的 21 个催化剂组合中，乙醇转化率从各自最低实测温度到最高实测温度均总体提高，其中 19 组逐点增加；$\mathrm{C}_4$ 烯烃选择性则表现出明显的组合差异，15 组逐点增加，其余组合出现局部回落或平台。以组内留一预测、原始散点和 0--100\% 边界共同筛选后，乙醇转化率有 15 组采用二次经验描述、5 组采用线性总体描述，A11 不强行确定具体低阶函数；选择性有 12 组采用二次描述、8 组采用线性总体描述，A6 仅作粗略描述。

附件2中，乙醇转化率由 43.5474\%（20 min）降至 29.9060\%（273 min），观测窗内的线性平均变化为 $X_{\mathrm{EtOH}}(t)=42.6709-0.05276t$，$R^2=0.9329$、$\mathrm{RMSE}_{\mathrm{LOO}}=2.025$ 个百分点。7 个离散时刻的 $\mathrm{C}_4$ 烯烃收率均依次下降；选择性在 36.72\%--40.32\% 间波动。因此，结论仅描述当前时间窗的变化，不将其外推为长期稳定性或催化剂失活机理。

\subsection{问题二：催化因素与温度的条件影响}

在 250--350 ℃的共同温区，21 个组合的乙醇转化率和 $\mathrm{C}_4$ 烯烃选择性端点差均为正，平均分别为 24.46 和 15.30 个百分点，温度端点促进方向在已有组合中最一致。严格匹配和跨背景对比同时显示，催化因素并不存在脱离背景的统一优劣：Co 负载量由 1 wt\% 增至 5 wt\% 时，转化率在两个匹配背景中分别出现 $+11.93$ 和 $-2.78$ 个百分点的平均差，而选择性在两个背景中的平均差均为负；乙醇进料条件和装料方式也均呈背景依赖。

在当前五因素定义下，仅 $(m_{\mathrm{cat}},q_{\mathrm{EtOH}})$ 与 $(q_{\mathrm{EtOH}},z)$ 形成完整局部四格。后者在 \SI{350}{\degreeCelsius} 的二阶对比为 $I_{X_{\mathrm{EtOH}}}=3.07$、$I_{S_{\mathrm{C4}}}=-19.41$ 个百分点，且两套四格删去任一共同温度后的平均二阶对比均保持原符号。上述结果描述的是离散水平、实际共同温度和指定背景下的局部非加性，不延伸为全局交互强度或严格因果判断。

\subsection{问题三：实验支持域内的收率优化}

以 $Y_{\mathrm{C4}}=X_{\mathrm{EtOH}}S_{\mathrm{C4}}/100$ 衡量联合条件的实测表现。无温度限制时，A3--\SI{400}{\degreeCelsius} 的实测 $\mathrm{C}_4$ 烯烃收率最高，为 44.72\%，对应 $X_{\mathrm{EtOH}}=83.70\%$、$S_{\mathrm{C4}}=53.43\%$；在各组合自身实测温区内的分段线性连续化下，该结论不产生新的内部峰值。

严格 $T<\SI{350}{\degreeCelsius}$ 时，实测最优为 A2--\SI{325}{\degreeCelsius}，收率 17.26\%。利用 A2 的 325--350 ℃相邻点构造的局部线性函数在 $T\to350^{-}$ 时给出 26.54\% 的上确界，但 \SI{350}{\degreeCelsius} 不属于严格可行域，故该数值不是可直接取得的最大收率。候选因素级代理在留出 A3 组时对其 44.72\% 的实际收率仅给出 17.99\%--26.80\% 的预测，因此不对未实验配方进行最优排序。

\subsection{问题四：5 次新增实验的安排}

第一阶段安排 4 次定向实验：A3--\SI{425}{\degreeCelsius} 和 A2--\SI{337.5}{\degreeCelsius} 用于分别细化高收益温区和严格低温区；A1--\SI{400}{\degreeCelsius} 与 A2--\SI{400}{\degreeCelsius} 用于补齐 Co 负载量和乙醇进料严格匹配中的关键缺格。完成后，按新增实测收率重新计算全温区与严格低温区的第一、第二名间隔；第 5 次独立复验间隔较小任务的当前第一名，并报告两次观测的均值和样本方差。该方案提供的是可执行的实验顺序，新增点的实际收率及最终复验对象须待实测后据实确定。
